\section{Related Work}
\label{sec:related}

Our work builds on evolutionary game theory, which has a long history.
\cite{fisher1930natural} explained the approximate equality of the sex ratio.
\cite{lewontin1961evolution} first applied game theory to evolution, and
\cite{smith1972game} together with \cite{smith1973logic} introduced the concept
of an Evolutionarily Stable Strategy (ESS).

\cite{bostrom2012} proposed the instrumental convergence thesis, which states that despite having
very different end goals, agents may still follow similar intermediate goals. This matches our
findings for potential-maximizing strategies as well. While we agree with Bostrom's
orthogonality theorem, that intelligence and goals are orthogonal, we also find that if we care
whether an intelligent strategy can actually survive in a competitive environment, this severely
restricts the space of goals that we can consider, and convergence may be inevitable.

\cite{hendrycks2023} argues
that AIs will likely be subject to natural selection (a key premise of our paper)
by showing that Lewontin's conditions apply to AI as well.

\cite{powerseeking} argues that reward-maximizing strategies are on average
drawn toward parts of the action space that offer more choices, on the
intuition that a larger action set is more likely to contain a high-reward
action. Because more actions confer more power, reward-maximizing strategies
tend to be power seeking. Although the argument differs from ours, both
approaches reach a similar conclusion.

\cite{Hutter:04uaibook}'s model of the universe and his use of horizons have
been a major inspiration for the definitions in this paper.

To our knowledge, nobody has applied Evolutionarily Stable Strategies to
problems relating to AI safety, as we have done in this paper.